\documentclass[11pt]{article}
\usepackage[margin=1in]{geometry}
\usepackage{amsmath,amssymb,amsthm,booktabs}
\usepackage[hidelinks]{hyperref}
\usepackage{xurl}
\setlength{\emergencystretch}{3em}
\newtheorem{theorem}{Theorem}
\newtheorem{corollary}[theorem]{Corollary}
\newcommand{\mr}{\operatorname{mr}}
\newcommand{\F}{\mathbb F}
\title{Fifty projective-cover certificates in a known MRC family}
\author{Exact values and graph-cover family audit}
\date{17 November 2026}
\begin{document}
\maketitle

Fifty additional named diagrams have exact meridional rank and bridge
number four. Their actual double branched covers are prime graph
manifolds. Consequently these equalities belong to the existing
Boileau--Jang--Weidmann graph-cover family, and are not presented as
cases outside known MRC families. The retained contribution is the
explicit value, finite projective-cover certificate and exact
family identification for each name. Worldwide priority of the
individual named computations is not established.

\begin{theorem}\label{thm:equalities}
For each knot in the following table,
\[
                         \mr(K)=b(K)=4.
\]
There is a transitive thirteen-sheet action of the knot group through
$\operatorname{SL}_3(\F_3)$ in which a meridian has cycle type
$2^4 1^5$. Its branched completion $Y$ has
$b_1(Y;\F_3)=2$.
\begin{center}\small
\begin{tabular}{lp{11.7cm}}
\toprule
table & indices\\
\midrule
$13a$ & $3167,3168,3169,3170,3188,3189,3190,3191,$\newline
$3293,3294,3295,3298,3347,3348,3351,3402,3403,3411$\\[3pt]
$13n$ & $326,329,330,341,344,345,348,351,352,355,357,358,$\newline
$510,513,514,518,520,521,524,525,526,529,530,531,$\newline
$3739,3742,3745,3748,3770,3773,3776,3779$\\
\bottomrule
\end{tabular}
\end{center}
\end{theorem}

\begin{theorem}\label{thm:graph}
For each of these fifty knots, the actual double branched cover
$\Sigma_2(K)$ is a closed orientable irreducible graph manifold.
A retained collection of three or four disjoint incompressible
normal tori cuts it into four or five explicitly recognized
Seifert fibre spaces.
\end{theorem}

\begin{corollary}\label{cor:known}
All fifty examples satisfy the prime graph-manifold double-cover
hypothesis of the published BJW18 theorem. With their upper bridge
bound four, the MRC equality is in that existing theorem's scope.
\end{corollary}

\section{Exact named diagrams and the bridge upper bound}

The finite input is \path{sl3_projective_family_inputs.json}.
Each row retains the literal thirteen-crossing PD, the independently
packaged Hoste--Thistlethwaite census PD and DT code, and an exact dart
bijection preserving edge pairing, cyclic crossing order and the
over/under distinction. Every selected bijection preserves orientation
and over/under. The verifier reads the current census package and checks
its pinned version and source-file hashes before verifying these maps.
A PD name is not inferred from a homology order or numerical invariant.

At each crossing the PD convention puts the overstrand in positions
one and three. Contracting the overstrand continuations gives thirteen
Wirtinger arcs. The replay checks planarity, connected projection,
one link component and every crossing incidence. Four seed arcs and
nine successive legal Wirtinger propagation steps reach all thirteen
arcs in each diagram. Thus each diagram has Wirtinger number at most
four. The published equivalence of Wirtinger number and bridge number,
BKVV20 Definition~1.2 and Theorem~1.3
\cite[printed pp.~245--246]{bkvv}, gives $b(K)\le4$.

\section{The thirteen-sheet homology certificate}

The projective points are the thirteen nonzero vectors of $\F_3^3$
modulo scalar, normalized by making the first nonzero coordinate one.
For every arc the input retains a literal $3$ by $3$ determinant-one
matrix and its permutation on these points. The replay checks the
projective action, nonidentity involution condition and every
Wirtinger conjugation relation directly. The four seed images generate
a permutation group of order $5616=|\operatorname{SL}_3(\F_3)|$.
Since the centre is trivial and the projective action is faithful,
this is the asserted group. Connectivity of the lifted generator
graph is checked separately by its spanning tree. Every meridian has
five fixed points and four transpositions, so $n=13$, $c=9$ and
$\delta=n-c=4$.

Use all thirteen diagram generators, without first eliminating the nine
nonseed generators. The full lifted graph has thirteen vertices and
$13\cdot13=169$ directed labelled edges. A checked spanning tree has
twelve edges, leaving $157$ columns. Each of the thirteen crossings
has thirteen lifts, giving $169$ relation rows. In addition, every cycle
of every generator is filled: there are $13\cdot9=117$ filling rows,
including all $13\cdot5=65$ fixed-sheet meridian rows.

The matrix is formed after these involutory meridian fillings.
Thus a crossing relator may be written $oao b$ with every letter
positive: the usual inverse letters differ by lifted meridian-square
relations, which are already consequences of the cycle fillings.
A negative directed-edge traversal and the corresponding positive
reverse traversal likewise differ by the same cycle relation.
This explains why the retained all-positive crossing walks and
all cycle fillings compute the homology of the ordinary branched
completion, rather than only a punctured cover or an orbifold cover.

The resulting $286$ by $157$ matrix has rank $155$ over $\F_3$ for
every row of the input. Hence $b_1(Y;\F_3)=157-155=2$.
The written general branched-cover bound in
\path{GENERAL_BRANCHED_COVER_RANK_BOUNDS.tex},
Theorem~1 (label \texttt{thm:general}), states
\[
 \mr(K)\ge
 \left\lceil\frac{n-1+b_1(Y;k)}{n-c}\right\rceil.
\]
Here it gives
\[
             \mr(K)\ge\left\lceil\frac{12+2}{4}\right\rceil=4.
\]
Together with $\mr(K)\le b(K)\le4$, this proves
Theorem~\ref{thm:equalities}.
No claim that the bare projective quotient needs four involutions is
used: this group does have generating triples of involutions.

\section{The actual double covers and their Seifert decompositions}

For each literal PD, Spherogram constructs the exterior with the diagram
meridian peripheral basis. Meridian coefficient-two orbifold filling
has a unique cyclic cover of degree two. Its branching slope has
primitive lift $(1,0)$; filling this lift constructs the actual
double branched cover. Regina converts and simplifies the resulting
closed triangulation. The retained input pins the random seed, exterior
isomorphism signature with peripheral data and closed-cover signature.
The replay reconstructs those signatures, rather than identifying
covers from matching homology. It checks that the result is valid,
closed, connected, orientable and irreducible.

Vertex normal surfaces are enumerated in quadrilateral coordinates.
The input supplies the chosen torus indices and complete quadrilateral
vectors. Each selected surface must be connected, orientable,
two-sided, have Euler characteristic zero and pass the exact
incompressibility predicate in its already checked closed irreducible
ambient manifold. All pairs must pass the disjointness test.

There is also a direct control on simultaneous representatives:
the sum normal surface must have components whose full triangle and
quadrilateral vectors are exactly those of the chosen surfaces,
as an unordered collection. Thus the replay cuts an actual disjoint
union of the specified tori. Pairwise coordinate compatibility alone
is not substituted for this test.

The sum is cut by \texttt{cutAlong()}, and every resulting connected
component is simplified and recognized by Regina as a Seifert fibre
space. Its exact finite triangulation signature, boundary count and
Seifert name must agree with the retained data. Every boundary component
is checked to be a torus. In particular no unrecognized component is
discarded. Three tori give four pieces; four give five. A redundant
$T^2\times I$ piece may occur, so minimality of this torus system or
identification with the canonical JSJ tori is not asserted or needed.

For example, $13a_{3167}$ has a sixteen-tetrahedron closed cover with
signature
\[
\texttt{qLLLLvALPQQkbcgfjinnlkklponpptsmgqsepbrjrsscnw}.
\]
Its tori have vertex-normal indices $1,15,32$.
The four recognized pieces are two copies of
$\operatorname{SFS}[D:(2,1)(3,-2)]$ and the two annulus-base spaces
$\operatorname{SFS}[A:(2,5)]$ and $\operatorname{SFS}[A:(2,1)]$.
All fifty cases have the corresponding exact cut-and-recognition
certificates.

A closed manifold obtained by gluing these Seifert spaces along their
torus boundaries is a graph manifold. The separately checked
irreducibility supplies primeness at the convention used by BJW.
This proves Theorem~\ref{thm:graph}.
BJW18 Theorem~1.2 and Corollary~1.3
\cite[printed p.~62]{bjw} apply to this graph-cover class.
Together with the standard meridional-rank one and two cases, their
rank-three statement gives MRC for a knot in this class with
$b(K)\le4$. This proves Corollary~\ref{cor:known}.
The homology calculation above determines the exact value four.

\section{Replay and assurance}

Run \path{verify_sl3_projective_family.py}; its output receipt is
\path{sl3_projective_family_verification.json}.
The pinned environment is SnapPy 3.3.2, Spherogram 2.4.1,
\texttt{snappy\_15\_knots} 1.2.1,
\texttt{snappy\_manifolds} 1.4 and Regina 7.4.
All fifty literal named diagrams, four-seed uppers, projective
representations, full branched-cover chains and actual graph-cover
decompositions are replayed. Wrong names, changed projective
permutations and corrupted normal vectors are rejected.

Trusted computational imports are the named census data, SnapPy's
combinatorial exterior/cover/filling construction and Regina's exact
normal-surface, irreducibility, disjointness and recognition algorithms.
The general-cover bound is a separate written proof; the BKVV
bridge equivalence and BJW family theorem are classical imports at
the scopes just stated. Numerical volume or numerical isometry is
not an acceptance test.

These are finite exact internal certificates with written proof
dependencies. They do not constitute formal verification, external
specialist review, worldwide named-priority verification, or a new
MRC family. No periodic-parameter generalization or exclusion of
other proof methods follows from these fifty cases.

\begin{thebibliography}{9}
\bibitem{bkvv}
R. Blair, A. Kjuchukova, R. Velazquez and P. Villanueva,
Wirtinger systems of generators of knot groups,
\emph{Communications in Analysis and Geometry} 28 (2020), 243--262,
Definition~1.2 and Theorem~1.3.
\url{https://doi.org/10.4310/CAG.2020.v28.n2.a2}.
\bibitem{bjw}
M. Boileau, Y. Jang and R. Weidmann,
Meridional rank and bridge number for a class of links,
\emph{Pacific Journal of Mathematics} 292 (2018), 61--84,
Theorem~1.2 and Corollary~1.3.
\url{https://doi.org/10.2140/pjm.2018.292.61}.
\end{thebibliography}
\end{document}
